\section{Results}
\label{sec:results}

% ── PLACEHOLDER ─────────────────────────────────────────────────────────────
% Complete once data collection is finished. Section order mirrors the RQs:
%   manipulation check (scoping) reported FIRST, then RQ1 → RQ2 → RQ3.

\subsection{Manipulation Check: Context Scoping Held}
\label{sec:results:scoping}

% Reported before any inferential result. From agent_context_sent snapshots:
% verify 100% of C1 turns contained no board/GitHub context and 100% of C2
% turns contained no GitHub context. This converts the IV from assumption to
% evidence.

\begin{table}[h]
  \caption{Scoping manipulation check: share of agent turns whose
    \texttt{agent\_context\_sent} snapshot contained each source, by condition.
    Cells outside the condition's allowlist must be 0\%. [PLACEHOLDER]}
  \label{tab:scoping}
  \begin{tabular}{lccc}
    \toprule
    Source in snapshot & \CI & \CII & \CIII \\
    \midrule
    Chat   & --- & --- & --- \\
    Board  & 0\% (expected) & --- & --- \\
    GitHub & 0\% (expected) & 0\% (expected) & --- \\
    \bottomrule
  \end{tabular}
\end{table}

\subsection{Participant Demographics and Data Quality}
\label{sec:results:participants}

% TODO: N per condition, attrition, task completion rates, data exclusions
%   (was_interrupted tasks, cross-device resumes), CONSORT-style flow.

\subsection{RQ1: Context Breadth and Effectiveness}
\label{sec:results:rq1}

% TODO: Primary linear mixed model on log baseline-normalized time
%   (condition * task_type). Lead with effect sizes + 95% CIs. The
%   condition × task-type interaction (H2) is the headline.
%   Correctness reported beside time (ordinal CLMM). Cognitive-load delta (H3).

\begin{figure}[h]
  \centering
  % TODO: Replace with actual figure — the H2 interaction plot.
  \fbox{\rule{0pt}{4cm}\rule{0.8\linewidth}{0pt}}
  \caption{Baseline-normalized task time by condition (\CI/\CII/\CIII) and task
    type (comprehension / single-surface / cross-surface). Error bars show
    95\% CI. The predicted condition~$\times$~task-type interaction
    concentrates on cross-surface tasks. [PLACEHOLDER]}
  \label{fig:time_by_condition}
\end{figure}

\begin{table}[h]
  \caption{Per-task correctness (0--3 rubric) by condition and task type, and
    within-person \NASATLX\ delta (agent~$-$~baseline) by condition. [PLACEHOLDER]}
  \label{tab:correctness_load}
  \begin{tabular}{lccc}
    \toprule
    Measure                       & \CI & \CII & \CIII \\
    \midrule
    Correctness, cross-surface (mean) & --- & --- & --- \\
    Correctness, single-surface (mean)& --- & --- & --- \\
    Correctness, comprehension (mean) & --- & --- & --- \\
    \midrule
    $\Delta$\NASATLX\ (agent $-$ baseline) & --- & --- & --- \\
    \bottomrule
  \end{tabular}
\end{table}

\subsection{RQ2: Reliance and Division of Labor}
\label{sec:results:rq2}

% TODO: Delegation ratio by condition; the agent funnel
%   (invoked → tool attempted → executed → accepted/edited/undone);
%   verification/override rates; whether reliance is calibrated to where the
%   agent is actually wrong (cross with correctness). Thematic analysis of
%   open-ended items (two coders, report κ).

\begin{table}[h]
  \caption{Reliance measures by condition: delegation ratio and the
    accept/edit/override breakdown of agent suggestions. [PLACEHOLDER]}
  \label{tab:reliance}
  \begin{tabular}{lccc}
    \toprule
    Measure                          & \CI & \CII & \CIII \\
    \midrule
    Delegation ratio (median)        & n/a & --- & --- \\
    Suggestions accepted (\%)        & --- & --- & --- \\
    Output edited before use (\%)    & --- & --- & --- \\
    Action undone / overridden (\%)  & --- & --- & --- \\
    \bottomrule
  \end{tabular}
\end{table}

\subsection{RQ3: Trust versus Reality}
\label{sec:results:rq3}

% TODO: Spearman correlations of trust composite and trust_speed_perceived
%   against measured normalized time and correctness. Characterize divergence
%   (H5): are agents that *do more* trusted even when no faster / occasionally
%   wrong? Pair int_switch_perceived with behavioral view-switch rate.

\begin{figure}[h]
  \centering
  \fbox{\rule{0pt}{4cm}\rule{0.8\linewidth}{0pt}}
  \caption{Perception--behavior gap: perceived speedup / AI-trust score (y) vs.\
    measured baseline-normalized time (x), per participant, colored by
    condition. Divergence from the diagonal is the H5 finding. [PLACEHOLDER]}
  \label{fig:perception_behavior}
\end{figure}

\begin{table}[h]
  \caption{AI-trust and integration-perception scale means (SD) by condition,
    with cross-surface telemetry counts for triangulation. [PLACEHOLDER]}
  \label{tab:trust}
  \begin{tabular}{lccc}
    \toprule
    Measure                          & \CI & \CII & \CIII \\
    \midrule
    AI trust scale (mean)            & --- & --- & --- \\
    \texttt{trust\_speed\_perceived} & --- & --- & --- \\
    Integration perception (mean)    & --- & --- & --- \\
    Cross-surface events (median)    & --- & --- & --- \\
    \bottomrule
  \end{tabular}
\end{table}
